

\section{Experiments}
\label{chap:results}

In this section, we present the empirical evaluation protocol used to assess L-GTA. It is designed to evaluate the method from three complementary perspectives: the controllability of latent space transformations, downstream forecasting utility, and the quality of the synthetic data. We also include an ablation study to highlight the relevance of our design choices in achieving these objectives.

\subsection{Experimental Setup}
\label{subsec:experimental_setup}

The experimental study is organized around the following research questions:
\begin{itemize}
    \item \textbf{Q1:} Do latent space manipulations in L-GTA produce synthetic series with controllable and interpretable transformation signatures?
    \item \textbf{Q2:} How does L-GTA compare with direct augmentation and SOTA augmentation methods in terms of downstream forecasting utility?
    \item \textbf{Q3:} How does L-GTA compare with direct augmentation and SOTA augmentation methods in terms of fidelity, diversity, privacy, and computational cost?
    \item \textbf{Q4:} Which components of L-GTA contribute most to controllability and reconstruction quality?
\end{itemize}

\subsubsection{Datasets and Baselines}

check TimeGEN format

\subsubsection{Evaluation Protocol}

Our primary benchmark is a downstream forecasting task. For each dataset, the series are partitioned into windows, and we reserve a number of them for evaluation purposes. Forecasting accuracy is measured with Mean Absolute Scaled Error (MASE) using an LSTM forecaster trained over multiple random restarts. We consider two evaluation protocols. In the \textbf{Train-on-Synthetic, Test-on-Real (TSTR)} setting, the forecaster is trained exclusively on synthetic windows and evaluated on held-out real windows. In the \textbf{augmentation} setting, the forecaster is trained on the union of original and synthetic windows and is again evaluated on held-out real windows.

To analyse the role of augmentation diversity, we evaluate both a \textbf{1-variant} and a \textbf{3-variant} setting. The 1-variant configuration applies a single latent transformation, while the 3-variant configuration combines jitter, scaling, and magnitude warping. We additionally evaluate L-GTA both with and without dynamic time features.

\subsection{Synthesis Quality Metrics}

Downstream forecasting alone does not fully characterise the quality of synthetic data. We therefore complement it with a feature-based synthesis-quality analysis. Each time series is embedded using time series features, after which we use to compare the synthetic and real datasets.

The resulting metrics are grouped into three categories:
\begin{itemize}
    \item \textbf{Fidelity:} improved precision, density, Frechet distance, Wasserstein distance, maximum mean discrepancy, and cosine similarity.
    \item \textbf{Diversity:} improved recall and coverage.
    \item \textbf{Privacy:} authenticity.
\end{itemize}

This analysis complements the downstream benchmark by quantifying whether a method produces synthetic series that are realistic, diverse, and sufficiently distinct from the original observations.

\subsection{Controllability and Ablation Studies}

To answer \textbf{Q1} and \textbf{Q4}, we perform two targeted analyses. The first evaluates whether the residual pattern induced by a latent transformation matches the intended direct transformation. Each transformed dataset is summarised by a three-dimensional fingerprint defined by residual autocorrelation, residual linearity, and amplitude dependence. These statistics are sufficient to distinguish additive perturbations such as jitter from smoother or multiplicative effects such as trend, drift, scaling, and magnitude warping. Therefore, this analysis tests whether L-GTA preserves the qualitative character of the transformation after decoding.

The second is a component ablation study over a family of L-GTA variants that differ in latent structure, equivariant loss, encoder design, channel attention, and the use of dynamic features. Transformation magnitude is also controlled, $\sigma \in \{0.5, 1.0, 2.0\}$, and for each variant we report four quantities: the mean Spearman correlation between transformation strength and observed effect, the proportion of monotonic responses, reconstruction MSE, and the fingerprint distance to the corresponding direct transformation.

\subsection{Practical Considerations}

Finally, we report computational cost in terms of training time.

\subsection{Replication Details}
\label{subsec:replication}

To facilitate replication, we summarise the main assumptions and hyperparameters used in our experiments.

\paragraph{L-GTA training protocol.}
Series are preprocessed with a sliding window of length $W = 10$ and scaled using a standard (zero mean, unit variance) scaler per series. The CVAE is trained with a latent dimension of $4$, batch size of $8$, and up to $1000$ epochs with early stopping. The optimizer is Adam with learning rate $0.001$ and weight decay $0.0001$. The learning rate is reduced by a factor of $0.5$ when training loss does not improve for $15$ epochs (minimum learning rate $10^{-6}$). Gradients are clipped to maximum norm $1.0$. The KL weight in the VAE loss is annealed linearly from $0$ to $0.1$ over the first $30$ epochs. After annealing, early stopping is reset and training continues until there are $30$ epochs without improvement in total loss. The equivariance regularisation weight is set to $1.0$ when equivariant training is used.

\paragraph{Downstream forecasting.}
The last $15\%$ of the windowed sequences are held out for evaluation; the remainder is used for training. The forecaster is a two-layer LSTM with $128$ hidden units per layer and a linear output head. It is trained with Adam (learning rate $10^{-3}$), batch size $32$, for $200$ epochs. Each configuration is evaluated over $5$ independent runs with a fixed random seed for data generation and training. 

\paragraph{Transformations and variants.}
In the 1-variant setting, a single transformation (e.g., jitter) is applied. In the 3-variant setting, we use jitter, scaling, and magnitude warping. In the ablation study we used $\sigma \in \{0.5, 1.0, 2.0\}$ and report metrics over $5$ repetitions per $\sigma$. 

\paragraph{Baselines and seeds.}
All methods used for comparison purposes (direct transformations, TimeGAN, TimeVAE and Diffusion-TS) use the same train/validation/test splits and the same random seed for reproducibility. L-GTA and the direct transformation share the same scaler and (when applicable) the same set of transformation types and $\sigma$ values for a fair comparison.
